\documentclass[a4paper, 10pt]{article}

% NALOGA 1
% Vključite pakete za podporo slovenščini,
% prepoznavo vhodnega kodiranja `utf8` in izhodnega kodiranja `T1`.
% Prekopirajte ustrezne tri vrstice iz pregledne datoteke `1-osnove~.tex`.

\usepackage{amsmath}
\usepackage[utf8]{inputenc}
\usepackage[T1]{fontenc}      % izhodna datoteka je kodirana v sistemu T1
\usepackage[slovene]{babel}
\usepackage{amsthm}
\usepackage{amsfonts}

\begin{document}

% Pri tej nalogi (in tudi vseh naslednjih) si pomagajte s paleto ukazov!
% Naslov, avtor, datum in povzetek naj bodo izdelani z uporabo ukazov in
% okolja, ki so v ta namen definirani v razredu `article`. Da se bo glava
% dokumenta s temi podatki prikazala, morate uporabiti ukaz `\maketitle`.

\title{
    ANALIZA I \\
     —\\
    \large IZREKI IN DEFINICIJE \\
}

\author{}
\date{}
\maketitle

{\theoremstyle{plain}
\newtheorem{izrek}{Izrek}
}


{\theoremstyle{definition}
\newtheorem{definicija}{Definicija}
}

\newpage
\section{LIMITA FUNKCIJE}

\begin{izrek}
    Naj bo f funkcija ene spremenljivke in a $\in \mathbb{R}$ stekališče $D_f$ z leve in desne. Velja: \\
    \[\exists\lim_{x \rightarrow a}f(x) \iff \exists\lim_{x \rightarrow a^-}f(x)  \wedge \exists\lim_{x \rightarrow a^+}f(x) \wedge \lim_{x \rightarrow a^-}f(x)  = \lim_{x \rightarrow a^+}f(x)  \] 
\end{izrek}

\begin{izrek}

\end{izrek}


\end{document}